\section*{Resumen}
\addcontentsline{toc}{section}{Resumen}

Este informe documenta el proceso metodológico del proyecto \textit{Desarrollo Social Económico}, que busca estudiar la asociación territorial entre contratación pública (SECOP II) y condiciones socioeconómicas en Colombia. La pregunta guía es: \textbf{¿Cómo se asocia la distribución territorial de la contratación pública con las condiciones socioeconómicas de las unidades territoriales colombianas?}

El informe cubre la selección de fuentes (CNPV 2018, IPM 2018 y DIVIPOLA), la homologación territorial, el tratamiento de datos faltantes, la construcción y control de indicadores y los diagnósticos preliminares. El universo contiene 1.122 códigos territoriales, incluidas áreas no municipalizadas; por ello no se describe como 1.122 municipios. Se revisaron nueve cuadros operativos. La especificación vigente registra 17 definiciones de indicadores; siete forman una matriz candidata provisional y las demás son descriptivas, contextuales o requieren revisión. Las fórmulas, denominadores y escenarios de no respuesta se documentan en los archivos reproducibles.

En la matriz base total, los diagnósticos reportan KMO $=0,642$, Bartlett $\chi^2=4664,873$ (21 gl; $p<0,001$), determinante $=0,0154$ y cero pares con $|r|\ge0,90$. El MSA de estrato 1--2 es $0,414$; 32 unidades territoriales presentan al menos un $|z|>4$. Son resultados preliminares: Bartlett puede ser sensible al tamaño del universo, un valor significativo no valida el PCA y el análisis de componentes aún no se ha ejecutado.

El Índice de Vulnerabilidad Socioeconómica (IVS) \textbf{no está construido}. Además, las variables censales miden condiciones estructurales transversales, no exposición a riesgos ni capacidad de respuesta; un índice futuro sería una aproximación parcial y deberá justificarse conceptualmente. El IPM se reserva para comparación externa, aunque la convergencia con él no sería una prueba independiente de validez. SECOP II y el análisis espacial siguen pendientes.

El informe distingue resultados reproducibles, decisiones provisionales y etapas futuras. Sus cifras deben interpretarse conforme a la fecha de corte y a las fuentes de evidencia indicadas.

\vspace{1cm}
\noindent\textbf{Fecha de corte:} 13 de octubre de 2026.\\
\textbf{Scripts de referencia:} \texttt{R/00\_config.R}, \texttt{R/01\_verificar\_consolidados.R}, \texttt{R/02\_base\_final\_indicadores.R}, \texttt{R/03\_auditoria\_indicadores.R}.\\
\textbf{Productos reproducibles:} \texttt{datos/trabajo/}, \texttt{resultados/base\_trabajo/}, \texttt{resultados/auditoria\_estadistica/}.
